\documentclass[12 pt, letterpaper]{hmcpset}
\usepackage{hw-preamble}
\setlist[enumerate, 1]{label = (\arabic*)}

\name{Jayden Thadani}
\class{MATH 147 --- Topology}
\assignment{Homework 3}
\duedate{13th February 2025}

\begin{document}

\begin{problem}[2.28]
	Let $A$ be a subset of a topological space $X$. Then $A^{\circ}$, $\partial A$ and $(X \setminus A)^{\circ}$ are disjoint sets whose union is $X$.
\end{problem}

\begin{solution}
	First we show that each pair of sets is disjoint.

	Notice that $A^{\circ}$ and $(X \setminus A)^{\circ}$ are disjoint since they are subsets of the disjoint sets $A$ and $X \setminus A$ respectively.

	Suppose $x \in \partial A = \overline{A} \cap \overline{X \setminus A}$. If $x$ is contained in either $A$ or $X \setminus A$ then $x$ isn't contained in the other, and thus, isn't contained in the interior of the other. Without loss of generality, say $x \in A$ and $x \notin X \setminus A$. Now $x \in \overline{X \setminus A}$ means $x$ is a limit point of $X \setminus A$ and every open neighbourhood of it contains points in $X \setminus A$. That is, no open neighbourhood of $x$ is contained in $A$ --- $x \notin A^{\circ}$.

	If $x \in \partial A$ is in neither $A$ nor $X \setminus A$, then it is a limit point of both. That is, each neighbourhood $U$ of $x$ contains points of both, $A$ and $X \setminus A$. Thus, no neighbourhood of $x$ is completely contained in $A$ or $X \setminus A$ and $x$ is not an interior point of $A$ nor $X \setminus A$. 

	Thus, $\partial A$ is disjoint from both $A^\circ$ and $(X \setminus A)^\circ$.

	Now we will show that the union of all of the sets is the whole space $X$. Specifically, we will show that if $x \in X$ is not in the interior of $A$ nor the interior of $X \setminus A$, then $x$ must be in $\partial A$.

	Suppose $x$ is not in the interior of $A$ nor in the interior of $X \setminus A$. Then there cannot be a neighbourhood $U$ of $x$ contained completely in $A$ or in $X \setminus A$. Thus, each neighbourhood $U$ of $x$ contains points of both $A$ and $X \setminus A$. $x$ is contained in exactly one of $A$ and $X \setminus A$ and (by the contrapositive of theorem 2.9) must be a limit point of the other. That is, $x \in \partial A = \overline{A} \cap \overline{X \setminus A}$.

	We have shown $X = A^{\circ} \cup (X \setminus A)^{\circ} \cup \partial A$.
\end{solution}

\pagebreak

\begin{problem}[3.3]
	Suppose $X$ is a set and $\mathcal{B}$ is a collection of subsets of $X$. Then $\mathcal{B}$ is a basis for some topology on $X$ if and only if
	\begin{enumerate}
		\item each point of $X$ is in some element of $\mathcal{B}$, and
		\item if $U$ and $V$ are sets in $\mathcal{B}$ and $p$ is a point in $U \cap V$, there is a set $W$ in $\mathcal{B}$ such that $p \in W \subseteq U \cap V$.
	\end{enumerate}
\end{problem}

\begin{solution}
	First we will show the conditions are necessary. Suppose the topology generated by $\mathcal{B}$ is $\mathcal{T}$. Since we must have $X \in \mathcal{T}$ equal to the union of basic open sets, we must indeed have every $x \in X$ contained in some basic open set. Since any $U, V \in \mathcal{B} \subseteq \mathcal{T}$ are open sets, so is their intersection, and thus, by theorem 3.1, $U \cap V$ is covered by basic open neighbourhoods $W \in \mathcal{B}$.

	Now suppose we have the conditions. We will show they are sufficient for $\mathcal{B}$ to generate a topology. Let $\mathcal{T}$ be the set of all unions of sets in $\mathcal{B}$. $X \in \mathcal{T}$ since each point $x \in X$ is in some basic open set, and $\emptyset \in \mathcal{T}$ as the empty union. 

	Since every open set is union of basic open sets, and thus, every union of these open sets is the union of the underlying basic open sets. That is, for an arbitrary collection of open sets $U_{\alpha} = \bigcup_{\beta \in \mathcal{J}} V_{\alpha, \beta}$ (indexed by $\alpha \in \mathcal{I}$) with $V_{\alpha, \beta} \in \mathcal{B}$ we have
	\[
		\bigcup_{\alpha \in \mathcal{I}} U_{\alpha} = \bigcup_{\alpha \in \mathcal{I}} \bigcup_{\beta \in \mathcal{J}} V_{\alpha, \beta} = \bigcup_{\alpha \in \mathcal{I}, \beta \in \mathcal{J}} V_{\alpha, \beta}
	\]
	which is open.

	Finally, if $U$ and $V$ are open, they are covered by basic open sets. Thus, each $x \in U \cap V$ is in basic open sets $U_{x} \subseteq U$ and $V_{x} \subseteq V$. Now, $x$ is in the set $U_{x} \cap V_{x} \subseteq U \cap V$. By the second condition, $U_{x} \cap V_{x}$ has a basic open neighbourhood $W_{x} \subseteq U_{x} \cap V_{x} \subseteq U \cap V$ containing $x$. The union of all these basic open neighbourhoods $W_{x}$ is just $U \cap V$, and thus, $U \cap V$ is open.
\end{solution}

\pagebreak

\begin{problem}[3.8]
	In the double-headed snake, show that every point is a closed set; however, it is impossible to find disjoint open sets $U$ and $V$ such that $0' \in U$ and $0'' \in V$.
\end{problem}

\begin{solution}
	We write $[0', b)$ to indicate $(0, b) \cup \{ 0' \}$ and similarly for $0''$.

	We will first show any $\{ a \}$ is closed in $\mathbb{R}_{+00}$. Suppose $b \neq a$. Suppose $b$ is one of the zeroes. If $a$ is not one of the zeroes $[b, a) \cap \{ a \}$ is empty, and thus, $b$ is not a limit point of $\{ a \}$. If $a$ is one of the zeroes, $[b, a) \cup \{ a \}$ is empty and $b$ is again not a limit point of $a$.

	Suppose $b$ is real. Suppose also that $a$ is real. If $b > a$ then $(a, b + 1) \cap \{ a \}$ is empty and if $b < a$ then $(0, b) \cap \{ a \}$ is empty, and thus, $b$ is not a limit point of $\{ a \}$. Thus, $\{ a \}$ contains all its limit points and is closed. If $a$ is one of the zeroes then $(0, b) \cup \{ a \}$ is empty and $b$ is again not a limit point of $a$.

	Suppose $U, V$ are open sets with $0' \in U$ and $0'' \in V$. Thus, for some $b_{1}, b_{2} \in \mathbb{R}$, we have $[0', b_{1}) \subseteq U$ and $[0'', b_{2}) \subseteq V$. Choosing any $a < \min \{ b_{1}, b_{2} \}$ we get $a$ in both basic open intervals, and thus, $a \in U \cap V$. That is, $U$ and $V$ cannot be disjoint.
\end{solution}

\pagebreak

\begin{problem}[3.9]
	\begin{enumerate}
		\item In the topological space $\mathbb{R}_{\text{har}}$, what is the closure of the set $H = \{ 1/n \}_{n \in \mathbb{N}}$?
		\item In the topological space $\mathbb{R}_{\text{har}}$, what is the closure of the set $-H = \{ -1/n \}_{n \in \mathbb{N}}$?
		\item Is it possible to find disjoint open sets $U$ and $V$ in $\mathbb{R}_{\text{har}}$ such that $0 \in U$ and $H \subseteq V$.
	\end{enumerate}
\end{problem}

\begin{solution}
	\begin{enumerate}
		\item
			\[
				\boxed{
					\overline{H} = H.
				}
			\]

			Since every $a \notin H$ has an open neighbourhood $(a - r, a + r) \setminus H$ which clearly has empty intersection with $H$, every $a \notin H$ is not a limit point of $H$.
		\item
			\[
				\boxed{
					\overline{- H} = -H \cup \{ 0 \}
				}
			\]

			Every open neighbourhood of $0$ in $\mathbb{R}_{\text{har}}$ contains $(-r, r) \setminus H$ for some $r > 0$. By the Archimedean property of $\mathbb{R}$, there is some $-1 / n > -r$, and thus, every open neighbourhood of $0$ has non-empty intersection with $-H$.
		\item It is not possible to find such disjoint open sets.

			Every open set $U$ containing $0$ contains $(-r, r) \setminus H$ for some $r > 0$. Again, by the Archimedean property, there is some $1 / n < r$. If $V$ is open and contains $H$, there is some $r' > 0$ such that $(1/n - r', 1/n + r') \setminus H$ is contained in $V$.

			The neighbourhoods $(-r, r) \setminus H$ and $(1/n - r', 1/n + r') \setminus H$ have non-empty intersection. For example, there is an irrational number between $1 / n$ and $\min \{ 1/n + r', r \}$, and that number cannot be in $H$ either.
	\end{enumerate}
\end{solution}

\pagebreak

\begin{problem}[3.10 (4)]
	In $\mathbb{H}_{\text{bub}}$, let $A$ and $B$ be countable sets on the $x$-axis such that $A$ and $B$ are disjoint. Show that there exist disjoint open sets $U$ and $V$ such that $A \subseteq U$ and $B \subseteq V$.
\end{problem}

\begin{solution}
	Consider the enumerations of $A = \{ a_{n} \}_{n \in \mathbb{N}}$ and $B = \{ b_{n} \}_{n \in \mathbb{N}}$. For each $n$ choose a sticky bubble $U_{a_{n}}$ stuck on $a_{n}$ that doesn't intersect with the sticky bubbles on any of the sticky bubbles $V_{b_{1}}, \dots, V_{b_{n - 1}}$ and choose a sticky bubble $V_{b_{n}}$ stuck on $b_{n}$ that doesn't intersect with any of the sticky bubbles $U_{a_{1}}, \dots, U_{a_{n}}$. 

	The construction in the previous paragraph is in fact possible. In each case, since there are finitely many previous sticky bubbles (in the Euclidean topology) the union of all of the closures of the sticky bubbles is closed and bounded --- compact, and thus has a finite distance $R$ from the point $b_{n}$ or $a_{n}$. By the triangle inequality, any sticky bubble of radius $r < R / 2$ does not intersect with any of the previous bubbles.

	Since each intersection $U_{a_{n}} \cap V_{b_{m}}$ is empty (by our construction), the intersection $\bigcup_{n \in \mathbb{N}} U_{a_{n}} \cap \bigcup_{n \in \mathbb{N}} V_{b_{n}}$ is also empty. That is, we have disjoint open sets $U = \bigcup_{n \in \mathbb{N}} U_{a_{n}}$ and $V = \bigcup_{n \in \mathbb{N}} V_{b_{n}}$ containing $A$ and $B$ respectively, as desired.
\end{solution}

\end{document}
